% Seção 3 — Referencial teórico e empírico. Reescrita em 2026-11-15: referencial
% de instrumento reduzido a um parágrafo; Calzada et al. corrigido pelo texto
% aceito (doc de 2026-11-10, §4). Magnitudes só da literatura.

\chapter{Referencial teórico e empírico}
\label{cap:referencial}

\section{A deriva como externalidade negativa}

O agrotóxico que deriva sobre o vizinho é um caso clássico de externalidade
negativa: o produtor que contrata a aplicação aérea internaliza o custo do
serviço e do produto, mas não o custo em saúde suportado por quem mora a favor
do vento. A solução negociada de \citeonline{coase1960} supõe custos de
transação baixos; aqui a população atingida é difusa e o dano é
probabilístico e de difícil atribuição, o que justifica a intervenção
regulatória. Qual instrumento usar, se proibição, taxa ou regulação de
desempenho, depende das inclinações relativas do benefício e do custo
marginais \cite{weitzman1974,skevas2012} e tem dimensão distributiva
\cite{lichtenberg1988,helfand1995}. Essa escolha é objeto de um
segundo ensaio da dissertação. O que este projeto fornece a ela é a premissa
empírica: saber se a proibição produziu o benefício em saúde que a motivou.

\section{Agrotóxicos e saúde ao nascer}

A literatura econômica recente sobre exposição gestacional a agrotóxicos
converge em três regularidades úteis ao desenho.

A primeira é a ordem de grandeza e a sua dependência da proximidade.
\citeonline{reynier2025glyphosate} estimam reduções médias de peso ao nascer da
ordem de 23 a 32 gramas associadas à difusão do glifosato em áreas rurais dos
Estados Unidos. \citeonline{calzada2023bananas} encontram déficit de cerca de
80 gramas, mas apenas para mães que moravam a até 150 metros do perímetro de
um bananal e tiveram três meses ou mais de fumigação intensiva durante a
gestação; para as que moravam mais longe não há gradiente. No Brasil,
\citeonline{rigotto2013trends} comparam municípios da fruticultura irrigada do
Ceará com municípios de agricultura familiar e documentam piores indicadores
de saúde nos primeiros, em estudo ecológico, sem pretensão causal.

A segunda é a concentração. Os efeitos aparecem no topo da distribuição de
exposição \cite{larsen2017agricultural} e na cauda inferior da distribuição de
peso esperado \cite{reynier2025glyphosate}; entre os poucos estudos que isolam
fungicidas está o de \citeonline{jones2020invasive}. A consequência para este
projeto é direta: se o efeito existe só para quem mora junto ao bananal, a
média de um município inteiro o dilui pela fração de nascimentos exposta, e um
desenho municipal pode não encontrá-lo mesmo que ele seja grande.

A terceira é o calendário. Quando a aplicação tem estação, o efeito é
identificado pelo cruzamento entre o período sensível da gestação e a estação
de aplicação: \citeonline{calzada2023bananas} encontram efeito nos dois
primeiros trimestres e estimativa sem significância estatística no terceiro, com a intensidade da fumigação
ditada pela estação chuvosa, e \citeonline{brainerd2014seasonal} usam o mês da
concepção e a estação de aplicação de agroquímicos na Índia. A ameaça
conhecida é a composição sazonal das mães: quem concebe em cada época do ano
difere em características que afetam o desfecho
\cite{buckles2013season,currie2013within}.

\section{Avaliação de uma política estadual e simultânea}

Como a lei vale para todo o estado desde a mesma data, os estimadores e as
decomposições de diferenças-em-diferenças para adoção escalonada
\cite{callaway2021multiple,goodmanbacon2021,dechaisemartin2020} não se aplicam:
não há variação de \emph{timing}. Restam a variação de \emph{calendário} entre
coortes de nascimento, para a qual a diferença tripla exige uma única hipótese
de tendências paralelas, a da diferença entre os dois contrastes que a compõem
\cite{olden2022triple}, e a variação de \emph{intensidade} prévia entre
municípios, para a qual \citeonline{callaway2024continuous} mostram que o
efeito no nível da dose recebida sai sob tendências paralelas usuais, ao passo
que a curva dose-resposta exige hipótese mais forte.

Quando poucos grupos são tratados, a inferência assintótica usual falha e são
necessários procedimentos próprios \cite{conley2011,ferman2019inference,
chernozhukov2021conformal,hagemann2025single}.
